\subsection{The Remainder Theorem} 
We begin with a theorem that will help us to get to grips with polynomials.
\begin{theorem}[Remainder Theorem]
For any polynomial $P(x)$ and number $k$, there is a polynomial $Q(x)$ and a number $R$ such that
\[
P(x)=Q(x)(x-k)+R
\]
and $R=P(k)$.
\end{theorem}
\begin{proof}
We can find $Q(x)$ and $R$ by dividing $P(x)\div\makeblank[1in]$. $Q(x)$ is the quotient, and $R$ is the remainder.

Remember that the degree of the remainder must be less than the degree of the divisor. \makeblank[1in] has degree \makeblank[.5in], so when we perform the division the degree of the remainder must be \makeblank[.5in]---that is, the remainder must be a \makeblank[1in]!

To see that $R=P(k)$, we just plug in $x=k$:
\begin{align*}
P(k)&=Q(k)(k-k)+R\\
&=Q(k)\cdot 0+R\\
&=0+R\\
&=R
\end{align*}
So we can find $Q(x)$ and $R$, as described by the theorem, for any polynomial.
\end{proof}
\vfill
\noindent... so what?